\documentclass[11pt,a4paper]{article}
\input{../../sablony/pripravy/priprava.sty}
\usepackage{tikz,pgfplots}
\pgfplotsset{compat=1.18}
\nastavPripravu{1. ročník}{Kvadratické rovnice a nerovnice}{106}{Kvadratická rovnice: obtížnější úpravy}
\begin{document}
\blok{Kvadratická rovnice: obtížnější úpravy}
Nejdříve upravit rovnici na nulový tvar a až potom určit $a,b,c$.
\blok{Metoda a návaznost na graf}
\begin{itemize}
\item Členy s $x^2$ musí skončit na jedné straně; pozor na znaménka při odečítání závorek.
\item Před použitím vzorce vytkni případný společný číselný násobek; kontroluj zkouškou.
\end{itemize}
\blok{Společně řešené příklady -- od jednoduchých k náročným}
Nejdříve zopakuj vzorec z minulé hodiny, pak přidej záporný koeficient $a$,
členy na obou stranách a nakonec složené mocniny.

\begin{prikladReseniBox}[title={Příklad 1}]
$2x^{2} - 7x + 3=0$
\tcblower
\[
\begin{aligned}
D&=49-4\cdot(2)\cdot(3)=25\\[2pt]
x_1&=\frac{1}{2}\\[2pt]
x_2&=3\\[2pt]
K&=\{\frac{1}{2};3\}
\end{aligned}
\]
\end{prikladReseniBox}
\begin{prikladReseniBox}[title={Příklad 2}]
$- 3x^{2} + 2x + 1=0$
\tcblower
\[
\begin{aligned}
D&=4-4\cdot(-3)\cdot(1)=16\\[2pt]
x_1&=- \frac{1}{3}\\[2pt]
x_2&=1\\[2pt]
K&=\{- \frac{1}{3};1\}
\end{aligned}
\]
\end{prikladReseniBox}
\begin{prikladReseniBox}[title={Příklad 3}]
$4x^{2} + 4x + 5=0$
\tcblower
\[
\begin{aligned}
D&=16-4\cdot(4)\cdot(5)=-64\\[2pt]
D&<0\ \Rightarrow\ \text{žádný reálný kořen}\\[2pt]
K&=\varnothing
\end{aligned}
\]
\end{prikladReseniBox}
\begin{prikladReseniBox}[title={Příklad 4 – neznámá na obou stranách}]
Vyřešte $(x-3)(x+2)=2(x-1)$.
\tcblower
\[\begin{aligned}
x^2-x-6&=2x-2\\x^2-3x-4&=0\\D&=9+16=25\\x_{1,2}&=\frac{3\pm5}{2}\\K&=\{-1;4\}
\end{aligned}\]
\end{prikladReseniBox}
\begin{prikladReseniBox}[title={Příklad 5 – složená rovnice se čtverci (dobrovolně)}]
Vyřešte $\displaystyle\frac12(2x-1)^2-\left[\frac12(x+1)\right]^2=3\left[\left(\frac x2\right)^2-\frac14\right]$.
\tcblower
\[\begin{aligned}
2x^2-2x+\frac12-\frac14(x^2+2x+1)&=\frac34(x^2-1)\\\frac74x^2-\frac52x+\frac14&=\frac34x^2-\frac34\\x^2-\frac52x+1&=0\\2x^2-5x+2&=0\\(2x-1)(x-2)&=0\\K&=\left\{\frac12;2\right\}
\end{aligned}\]
\end{prikladReseniBox}
\poznamka{Vynech Příklad 5, pokud nezbude čas; je to přirozený vrchol stupňování obtížnosti.}
\blok{Samostatné procvičování}
\begin{zadaniReseniBox}[title={Úloha 1}]
$3x^{2} - 5x - 2=0$
\tcblower
\[
\begin{aligned}
D&=25-4\cdot(3)\cdot(-2)=49\\[2pt]
x_1&=- \frac{1}{3}\\[2pt]
x_2&=2\\[2pt]
K&=\{- \frac{1}{3};2\}
\end{aligned}
\]
\end{zadaniReseniBox}
\begin{zadaniReseniBox}[title={Úloha 2}]
$2x^{2} + 5x - 3=0$
\tcblower
\[
\begin{aligned}
D&=25-4\cdot(2)\cdot(-3)=49\\[2pt]
x_1&=-3\\[2pt]
x_2&=\frac{1}{2}\\[2pt]
K&=\{-3;\frac{1}{2}\}
\end{aligned}
\]
\end{zadaniReseniBox}
\blok{Volitelná rezerva}
\poznamka{Studenti si vymyslí obdobný příklad se změnou jednoho koeficientu a vysvětlí, jak se změní počet kořenů či řešení.}
\poznamka{Další typové úlohy jsou v mé učebnici Work-MAT: kapitola \emph{Kvadratické rovnice}. V repozitáři ji najdeš pod \emph{Kvadratické funkce a rovnice}; navazují také úlohy A007--A010 ze sbírky.}
\blok{Domácí úkol}
\begin{zadaniReseniBox}[title={Domácí příprava}]
\begin{enumerate}[label=\alph*),itemsep=2pt]
\item $3x^{2} - 10x + 3=0$
\item $6x^{2} + x - 1=0$
\item $- 2x^{2} + 5x + 3=0$
\item $4x^{2} - 12x + 9=0$
\item $5x^{2} + 4x + 3=0$
\item $3x^{2} - 15=0$
\item Řešte v $\mathbb R$: $4x^{2} - 4x + 1=0$.

\item $5x-3x^2=2$.
\item $(x-4)(x+1)=3(x-2)$.
\item $2(x-1)^2-(x+2)^2=3x-10$.
\end{enumerate}
\tcblower
\begin{enumerate}[label=\alph*),itemsep=2pt]
\item
\[\begin{aligned}
D&=100-4\cdot(3)\cdot(3)=64\\[2pt]
x_1&=\frac{1}{3}\\[2pt]
x_2&=3\\[2pt]
K&=\{\frac{1}{3};3\}
\end{aligned}\]
\item
\[\begin{aligned}
D&=1-4\cdot(6)\cdot(-1)=25\\[2pt]
x_1&=- \frac{1}{2}\\[2pt]
x_2&=\frac{1}{3}\\[2pt]
K&=\{- \frac{1}{2};\frac{1}{3}\}
\end{aligned}\]
\item
\[\begin{aligned}
D&=25-4\cdot(-2)\cdot(3)=49\\[2pt]
x_1&=- \frac{1}{2}\\[2pt]
x_2&=3\\[2pt]
K&=\{- \frac{1}{2};3\}
\end{aligned}\]
\item
\[\begin{aligned}
D&=144-4\cdot(4)\cdot(9)=0\\[2pt]
x&=\frac{3}{2}\\[2pt]
K&=\{\frac{3}{2}\}
\end{aligned}\]
\item
\[\begin{aligned}
D&=16-4\cdot(5)\cdot(3)=-44\\[2pt]
D&<0\ \Rightarrow\ \text{žádný reálný kořen}\\[2pt]
K&=\varnothing
\end{aligned}\]
\item
\[\begin{aligned}
x^2&=5\\|x|&=\sqrt{5}\\[2pt]
K&=\{- \sqrt{5};\sqrt{5}\}
\end{aligned}\]
\item \[ \begin{aligned}D&=(-4)^2-4\cdot(4)\cdot(1)=0\\K&=\{\frac{1}{2}\}\end{aligned}\]

\item \[\begin{aligned}3x^2-5x+2&=0\\D&=25-24=1\\x_{1,2}&=\frac{5\pm1}{6}\\K&=\left\{\frac23;1\right\}\end{aligned}\]
\item \[\begin{aligned}x^2-3x-4&=3x-6\\x^2-6x+2&=0\\D&=36-8=28\\x_{1,2}&=3\pm\sqrt7\\K&=\{3-\sqrt7;3+\sqrt7\}\end{aligned}\]
\item \[\begin{aligned}2(x^2-2x+1)-(x^2+4x+4)&=3x-10\\x^2-8x-2&=3x-10\\x^2-11x+8&=0\\D&=121-32=89\\K&=\left\{\frac{11-\sqrt{89}}2;\frac{11+\sqrt{89}}2\right\}\end{aligned}\]
\end{enumerate}
\end{zadaniReseniBox}
\poznamka{Orientační čas: úvod 5 min, společné příklady 18 min, samostatně 15 min, kontrola a shrnutí 7 min. Rezerva je dobrovolná. Číselné parametry modelových aplikací nejsou skutečně naměřená data.}
\end{document}
